\section{Large Language Models and Prompt Engineering in Healthcare}\label{sec:rw_llm}

The third pillar of the proposed platform concerns the translation of statistical output into clinician-readable language and the translation of natural-language questions into safe database queries. Both rely on large language models, whose behavior in clinical settings is critically dependent on prompt design. The work reviewed here provides the most rigorous recent treatment of that dependency.

Liu et al. \cite{liu2025} present a tutorial framework for prompt engineering in clinical practice, aimed at clinicians who are not experts in natural language processing. The authors observe that the effectiveness of LLMs such as the GPT series and Google PaLM hinges on the quality of the prompt, since small changes in wording or structure can substantially alter output quality, which is a particularly dangerous property in a medical context.

The framework consolidates a range of prompting techniques, including zero-shot, one-shot, and few-shot prompting, chain-of-thought, self-consistency, generated-knowledge, and meta-prompting, and provides actionable guidance on defining objectives, applying core principles, iteratively refining prompts, and integrating LLMs into interoperable EHR systems. A central concern of the tutorial is the mitigation of hallucination and the preservation of ethical standards and patient safety, which is exactly the risk class that HospitAI must control when an LLM is placed in front of a clinical database.

Although the framework proposed by Liu et al.\ \cite{liu2025} is essential for direct textual clinical interactions, unlike the solution proposed in this thesis (HospitAI) it does not cover the interpretation of structured data produced by unsupervised machine learning algorithms. HospitAI extends this framework by applying it directly to the centroids extracted through K-Prototypes and the sequences produced by PrefixSpan, creating an automated bridge between data mining and constraint-based clinical interpretation. Furthermore, the tutorial treats the LLM as a free-form text generator, whereas HospitAI encloses the stochastic model inside a deterministic control graph in which every generated SQL query is validated at the abstract-syntax-tree level for read-only safety before execution, with a bounded retry loop on failure, as formalized in Chapter~\ref{ch:analysis} and implemented in Chapter~\ref{ch:design}.
